\documentclass[../../main.tex]{subfiles}

\begin{document}

\section{整数与实数}\label{section:ch01-整数与实数}

迄今我们所讨论的可以说是拓扑学的\textbf{逻辑基础}——集合论的一些基本概念。现在转入讨论它的\textbf{数学基础}——整数与实数系。相关的内容在前几节的例题和习题中已经非正式地使用过，现在需要给予正式处理。

建立这些基础的一个办法是，仅仅应用集合论公理来构造实数系，也就是说，赤手空拳地干。这样处理问题要花费很多时间和精力，并且其中逻辑的味道远远超过数学。

第二种方法比较简单，它假定已有关于实数的一些公理，从这些公理出发进行讨论。本节大体上就是这样处理实数的。准确些说，我们将给出关于实数的一些公理，并且指出如何从这些公理出发导出整数和实数的一些熟知的性质。但是，大部分证明将留作习题。如果你以前已经见过相关的内容，那么我们的讨论将会加深你的理解。如果你不知道，那么把这些习题详细地做一遍，将会加强你在数学基础方面的知识。

首先，我们将要用到集合论中的一个定义。

\begin{definition}[二元运算]\label{definition:ch01-二元运算}
集合 $A$ 中的一个二元运算（binary operation）是将 $A \times A$ 映到 $A$ 中的一个函数 $f$。
\end{definition}

对于集合 $A$ 中的二元运算 $f$，往往使用一种与\cref{section:ch01-函数}引进的标准函数记号不同的记号。函数 $f$ 在点 $(a, a')$ 处的值不用 $f(a, a')$，而是把函数符号写在点的两个坐标中间，即用 $afa'$ 来表示运算的值。与关系的情形一样，今后也经常使用一些不同于字母的符号来表示一个运算。常用的符号有加号（$+$）、乘号（$\cdot$）和（$\circ$），以及星号（$\ast$）等等。

\subsection*{假定}

我们假定有一个称之为实数（real numbers）的集合 $\mathbb{R}$，在 $\mathbb{R}$ 上有分别称之为加法运算和乘法运算的两个二元运算 $+$ 和 $\cdot$，以及 $\mathbb{R}$ 上的一个全序关系 $<$，它们具有以下性质：

\textbf{代数性质}
\begin{enumerate}[(1)]
\item $(x + y) + z = x + (y + z)$ 和 $(x \cdot y) \cdot z = x \cdot (y \cdot z)$ 对于 $\mathbb{R}$ 中的所有 $x$，$y$，$z$ 成立.
\item $x + y = y + x$ 和 $x \cdot y = y \cdot x$ 对于 $\mathbb{R}$ 中的所有 $x$，$y$ 成立.
\item $\mathbb{R}$ 中有唯一的一个元素，称为零（zero），记作 $0$，它使得对于所有 $x \in \mathbb{R}$，$x + 0 = x$.

$\mathbb{R}$ 中有唯一一个不是 $0$ 的元素，称为一（one），记作 $1$，它使得对于所有 $x \in \mathbb{R}$，$x \cdot 1 = x$.
\item 对于 $\mathbb{R}$ 中每一个 $x$，存在 $\mathbb{R}$ 中唯一的一个 $y$，使得 $x + y = 0$；对于 $\mathbb{R}$ 中每一个 $x \neq 0$，存在 $\mathbb{R}$ 中唯一的一个 $y$，使得 $x \cdot y = 1$.
\item $x \cdot (y + z) = (x \cdot y) + (x \cdot z)$ 对于所有 $x$，$y$，$z \in \mathbb{R}$ 成立.
\end{enumerate}

\textbf{代数与序的混合性质}
\begin{enumerate}[(6)]
\item 若 $x > y$，则 $x + z > y + z$.

若 $x > y$ 和 $z > 0$，则 $x \cdot z > y \cdot z$.
\end{enumerate}

\textbf{序性质}
\begin{enumerate}[(7)]
\item 全序关系 $<$ 具有上确界性质.
\item 若 $x < y$，则存在一个元素 $z$，使得 $x < z$ 和 $z < y$.
\end{enumerate}

根据性质 (1)$\sim$(5) 即可得到熟知的“代数定律”。给定 $x$，用 $-x$ 记满足 $x + y = 0$ 的数 $y$。$-x$ 称为 $x$ 的负元（negative element）。我们通过公式 $z - x = z + (-x)$ 定义减法运算（subtraction operation）。类似地，对于 $x \neq 0$，用 $1/x$ 记使得 $x \cdot y = 1$ 成立的数 $y$，$1/x$ 称为 $x$ 的倒数（reciprocal）。我们通过公式 $z/x = z \cdot (1/x)$ 定义商（quotient）。常用的符号法则以及分数的加法与乘法法则都可以作为定理推出来。这些代数定律写在本节末尾的习题 1 中。$x \cdot y$ 常简单地记为 $xy$。

将性质 (6) 与性质 (1)$\sim$(5) 合在一起，可以证明通常的“不等式法则”，比如
\begin{align*}
\text{若 } x > y \quad \text{和} \quad z < 0, \text{则 } x \cdot z < y \cdot z.
\end{align*}
\begin{align*}
-1 < 0 \quad \text{和} \quad 0 < 1.
\end{align*}
所有不等式法则写在习题 2 中。

若 $x > 0$，则定义 $x$ 为正数（positive number）；若 $x < 0$，则定义 $x$ 为负数（negative number）。全体正实数记作 $\mathbb{R}_{+}$，全体非负实数记作 $\bar{\mathbb{R}}_{+}$（理由将在后面说明）。性质 (1)$\sim$(6) 是近世代数中熟知的性质。具有两个满足性质 (1)$\sim$(5) 的二元运算的任何集合，在代数学中称为一个域（field），具有满足性质 (6) 的全序关系的域，称为一个有序域（ordered field）。

另一方面，性质 (7) 和 (8) 是拓扑学中熟知的性质，它们仅涉及全序关系，具有满足性质 (7) 和 (8) 的全序关系的任何集合，在拓扑学中称为线性连续统（linear continuum）。

将关于有序域的公理（性质 (1)$\sim$(6)）与关于线性连续统的公理（性质 (7) 和 (8)）合起来，发现有些结果重复了。特别地，性质 (8) 可以由另外几个性质推证出来；对于 $x < y$，能够证明 $z = (x + y)/(1 + 1)$ 满足性质 (8) 的条件。因此，在实数理论中，仅将性质 (1)$\sim$(7) 作为公理假设，而将性质 (8) 作为一个定理。我们之所以在实数的基本性质中还要提到性质 (8)，那是为了要强调：性质 (8) 和上确界性质是实数中全序关系的两个主要性质。从它们可以导出 $\mathbb{R}$ 的许多拓扑性质，在第 3 章中将会看到这一点。

在这一系列性质中，并没有告诉我们什么是整数。现在我们仅仅用性质 (1)$\sim$(6) 来定义整数。

\begin{definition}[归纳集与正整数]\label{definition:ch01-归纳集}
实数集的一个子集 $A$ 称为归纳的（inductive），如果它包含着数 $1$，并且只要 $x \in A$ 则必有 $x + 1 \in A$。设 $\mathcal{A}$ 为 $\mathbb{R}$ 中所有包含 $1$ 的归纳子集的族。正整数（positive integer）集 $\mathbb{Z}_{+}$ 定义为
\begin{align}
\mathbb{Z}_{+} = \bigcap_{A \in \mathcal{A}} A.\label{align:ch01-正整数定义}
\end{align}
\end{definition}

注意，正实数集 $\mathbb{R}_{+}$ 包含 $1$ 并且是归纳集（若 $x > 0$，则 $x + 1 > 0$），于是 $\mathbb{R}_{+}$ 属于 $\mathcal{A}$。因此 $\mathbb{Z}_{+} \subset \mathbb{R}_{+}$，$\mathbb{Z}_{+}$ 的元素都是正数，这正是我们选用这个术语的原因。因为所有实数 $x$（$x \geqslant 1$）的集合是归纳集并且包含 $1$，所以 $1$ 就是 $\mathbb{Z}_{+}$ 的最小元。

从这个定义出发不难得出 $\mathbb{Z}_{+}$ 的基本性质：
\begin{enumerate}[(1)]
\item $\mathbb{Z}_{+}$ 是归纳集.
\item（归纳原理）若 $A$ 是包含 $1$ 的正整数的一个归纳集，则 $A = \mathbb{Z}_{+}$.
\end{enumerate}

整数（integer）集 $\mathbb{Z}$ 定义为由正整数 $\mathbb{Z}_{+}$、数 $0$ 及 $\mathbb{Z}_{+}$ 中元素的负数组成的集合。可以证明两个整数的和、差、积是整数，但其商未必是整数。整数的商的集合 $\mathbb{Q}$ 称为有理数（rational number）集。

可以证明：对于给定的整数 $n$，不存在满足 $n < a < n + 1$ 的整数 $a$。

若 $n$ 是一个正整数，我们以 $S_{n}$ 表示所有小于 $n$ 的正整数的集合，称作正整数的一个截（section）。那么，$S_{1}$ 为空集，$S_{n+1}$ 为由 $1$ 到 $n$ 的所有正整数所组成的集合。在后续讨论中，我们将使用记号
\begin{align*}
\{1, 2, \cdots, n\} = S_{n+1}.
\end{align*}

以下将要证明的两条性质大家未必很熟悉，但它们却非常有用。可以将其理解为归纳原理的另一种表现形式。

\begin{theorem}[良序性质]\label{theorem:ch01-良序性质}
$\mathbb{Z}_{+}$ 的每一个非空子集有一个最小元。
\end{theorem}
\begin{proof}

我们首先证明下述论断成立：对于每一个 $n \in \mathbb{Z}_{+}$，$\{1, \cdots, n\}$ 的每一个非空子集有最小元。

设 $A$ 是使上述论断成立的所有正整数 $n$ 的集合，则 $A$ 包含 $1$。这是由于 $n = 1$ 时，$\{1, \cdots, n\}$ 仅有的非空子集是集合 $\{1\}$ 自身。其次，设 $A$ 包含 $n$，我们证明它包含 $n + 1$。为此，设 $C$ 为 $\{1, \cdots, n + 1\}$ 的一个非空子集。如果 $C$ 由一个元素 $n + 1$ 组成，那么这个元素就是 $C$ 的最小元。如果 $C$ 多于一个元素，考虑非空集 $C \cap \{1, \cdots, n\}$。因为 $n \in A$，这个集合有一个最小元，它自然也就是 $C$ 的最小元。于是 $A$ 为包含 $1$ 的归纳集，所以 $A = \mathbb{Z}_{+}$，这就证明了对于所有 $n \in \mathbb{Z}_{+}$ 论断为真。

现在证明定理。设 $D$ 为 $\mathbb{Z}_{+}$ 的一个非空子集。在 $D$ 中取一个元素 $n$，则集合 $A = D \cap \{1, \cdots, n\}$ 非空，从而 $A$ 有一个最小元 $k$。$k$ 自然就是 $D$ 的最小元。
\end{proof}

\begin{theorem}[强归纳原理]\label{theorem:ch01-强归纳}
设 $A$ 是一个以正整数为元素的集合。假定对于每一正整数 $n$，$S_{n} \subset A$ 蕴涵 $n \in A$。则 $A = \mathbb{Z}_{+}$。
\end{theorem}

\begin{proof}
假定 $A$ 与 $\mathbb{Z}_{+}$ 不等，那么必定存在一个不属于 $A$ 的最小正整数，记作 $n$。由于每一小于 $n$ 的正整数属于 $A$，因此 $S_{n} \subset A$。但我们的假设蕴涵着 $n \in A$，矛盾。
\end{proof}

到现在为止，我们仅仅用到了关于有序域的公理——实数性质 (1)$\sim$(6)。什么时候用到上确界公理 (7) 呢？

可以用上确界公理证明正整数集 $\mathbb{Z}_{+}$ 在 $\mathbb{R}$ 中没有上界。这是实直线的 Archimedean 有序性质（Archimedean ordering property）。为了证明这一点，我们假定 $\mathbb{Z}_{+}$ 有一个上界，从而导致矛盾。如果 $\mathbb{Z}_{+}$ 有一个上界，则有上确界 $b$。于是存在 $n \in \mathbb{Z}_{+}$，使 $n > b - 1$；若不然，则 $b - 1$ 就是 $\mathbb{Z}_{+}$ 的一个小于 $b$ 的上界。因此 $n + 1 > b$，与 $b$ 是 $\mathbb{Z}_{+}$ 的上界矛盾。

上确界公理还可以用来证明关于 $\mathbb{R}$ 的另一些结果。例如用以证明，对于任意正实数 $x$，其正平方根 $\sqrt{x}$ 的存在唯一性。而这个事实又可以用来证明不是有理数的实数的存在性，一个很容易的例子就是数 $\sqrt{2}$。

我们用符号 $2$ 表示 $1 + 1$，符号 $3$ 表示 $2 + 1$ 等等，就得到正整数的标准符号。按照这个程序，每一个正整数都确定了一个唯一的记号，然而，我们永远用不到这一点，也不予以证明。

关于整数和实数的这些性质，以及后面要用到的一些其他性质，都将在习题中概括。

\subsection*{习题}

\begin{exercise}\label{exercise:ch01-整数与实数-1}
仅用公理 (1)$\sim$(5) 证明 $\mathbb{R}$ 的下述“代数定律”：
\begin{enumerate}[(a)]
\item 若 $x + y = x$，则 $y = 0$.
\item $0 \cdot x = 0$. [提示：计算 $(x + 0) \cdot x$.]
\item $-0 = 0$.
\item $-(-x) = x$.
\item $x(-y) = -(xy) = (-x)y$.
\item $(-1)x = -x$.
\item $x(y - z) = xy - xz$.
\item $-(x + y) = -x - y$，$-(x - y) = -x + y$.
\item 若 $x \neq 0$ 并且 $x \cdot y = x$，则 $y = 1$.
\item 若 $x \neq 0$，则 $x/x = 1$.
\item $x/1 = x$.
\item $x \neq 0$ 并且 $y \neq 0 \Longrightarrow xy \neq 0$.
\item 若 $y$，$z \neq 0$，则 $(1/y)(1/z) = 1/(yz)$.
\item 若 $y$，$z \neq 0$，则 $(x/y)(w/z) = (xw)/(yz)$.
\item 若 $y$，$z \neq 0$，则 $(x/y) + (w/z) = (xz + wy)/(yz)$.
\item $x \neq 0 \Longrightarrow 1/x \neq 0$.
\item 若 $w$，$z \neq 0$，则 $1/(w/z) = z/w$.
\item 若 $y$，$w$，$z \neq 0$，则 $(x/y)/(w/z) = (xz)/(yw)$.
\item 若 $y \neq 0$，则 $(ax)/y = a(x/y)$.
\item 若 $y \neq 0$，则 $(-x)/y = x/(-y) = -(x/y)$.
\end{enumerate}
\end{exercise}

\begin{exercise}\label{exercise:ch01-整数与实数-2}
应用公理 (1)$\sim$(6) 及习题\cref{exercise:ch01-整数与实数-1}中的那些结论，证明关于 $\mathbb{R}$ 的“不等式法则”：
\begin{enumerate}[(a)]
\item $x > y$ 和 $w > z \Longrightarrow x + w > y + z$.
\item $x > 0$ 和 $y > 0 \Longrightarrow x + y > 0$ 和 $x \cdot y > 0$.
\item $x > 0 \Longleftrightarrow -x < 0$.
\item $x > y \Longleftrightarrow -x < -y$.
\item $x > y$ 和 $z < 0 \Longrightarrow xz < yz$.
\item $x \neq 0 \Longrightarrow x^{2} > 0$，其中 $x^{2} = x \cdot x$.
\item $-1 < 0 < 1$.
\item $xy > 0 \Longleftrightarrow x$ 与 $y$ 同为正数或同为负数.
\item $x > 0 \Longrightarrow 1/x > 0$.
\item $x > y > 0 \Longrightarrow 1/x < 1/y$.
\item $x < y \Longrightarrow x < (x + y)/2 < y$.
\end{enumerate}
\end{exercise}

\begin{exercise}\label{exercise:ch01-整数与实数-3}
(a) 证明：若 $\mathcal{A}$ 是归纳集的一个族，则 $\mathcal{A}$ 中元素的交是一个归纳集。

(b) 证明 $\mathbb{Z}_{+}$ 的基本性质 (1) 和 (2)。
\end{exercise}

\begin{exercise}\label{exercise:ch01-整数与实数-4}
(a) 用归纳法证明：对于每一个 $n \in \mathbb{Z}_{+}$，$\{1, \cdots, n\}$ 的每一个非空子集有最大元。

(b) 说明为什么不能由 (a) 推出 $\mathbb{Z}_{+}$ 的每一个非空子集有最大元？
\end{exercise}

\begin{exercise}\label{exercise:ch01-整数与实数-5}
证明 $\mathbb{Z}$ 与 $\mathbb{Z}_{+}$ 的下述性质：
\begin{enumerate}[(a)]
\item $a, b \in \mathbb{Z}_{+} \Longrightarrow a + b \in \mathbb{Z}_{+}$. [提示：证明对于给定的 $a \in \mathbb{Z}_{+}$，集合 $X = \{x \mid x \in \mathbb{R}$ 和 $a + x \in \mathbb{Z}_{+}\}$ 是一个归纳集.]
\item $a, b \in \mathbb{Z}_{+} \Longrightarrow a \cdot b \in \mathbb{Z}_{+}$.
\item 证明：$a \in \mathbb{Z}_{+} \Longrightarrow a - 1 \in \mathbb{Z}_{+} \cup \{0\}$. [提示：令 $X = \{x \mid x \in \mathbb{R}$ 和 $x - 1 \in \mathbb{Z}_{+} \cup \{0\}\}$. 证明 $X$ 是一个归纳集.]
\item $c, d \in \mathbb{Z} \Longrightarrow c + d \in \mathbb{Z}$ 和 $c - d \in \mathbb{Z}$. [提示：首先对 $d = 1$ 证明结论.]
\item $c, d \in \mathbb{Z} \Longrightarrow c \cdot d \in \mathbb{Z}$.
\end{enumerate}
\end{exercise}

\begin{exercise}\label{exercise:ch01-整数与实数-6}
设 $a \in \mathbb{R}$。对于 $n \in \mathbb{Z}_{+}$ 归纳地定义
\begin{align*}
a^{1} = a,
\end{align*}
\begin{align*}
a^{n+1} = a^{n} \cdot a.
\end{align*}
（关于归纳定义过程的讨论，参见第 7 节。）证明：对于 $n$，$m \in \mathbb{Z}_{+}$ 和 $a$，$b \in \mathbb{R}$，
\begin{align*}
a^{n}a^{m} = a^{n+m},
\end{align*}
\begin{align*}
(a^{n})^{m} = a^{nm},
\end{align*}
\begin{align*}
a^{m}b^{m} = (ab)^{m}.
\end{align*}
这些称为指数法则（laws of exponents）。[提示：固定 $n$，对 $m$ 归纳证明上式。]
\end{exercise}

\begin{exercise}\label{exercise:ch01-整数与实数-7}
设 $a \in \mathbb{R}$ 并且 $a \neq 0$。定义 $a^{0} = 1$。对于 $n \in \mathbb{Z}_{+}$，定义 $a^{-n} = 1/a^{n}$。证明对于 $a$，$b \neq 0$ 以及 $n$，$m \in \mathbb{Z}$，指数法则成立。
\end{exercise}

\begin{exercise}\label{exercise:ch01-整数与实数-8}
(a) 证明 $\mathbb{R}$ 有下确界性质。

(b) 证明 $\inf\{1/n \mid n \in \mathbb{Z}_{+}\} = 0$。

(c) 证明：对于给定的 $a$（$0 < a < 1$），有 $\inf\{a^{n} \mid n \in \mathbb{Z}_{+}\} = 0$。[提示：令 $h = (1 - a)/a$，证明 $(1 + h)^{n} \geqslant 1 + nh$.]
\end{exercise}

\begin{exercise}\label{exercise:ch01-整数与实数-9}
(a) 证明：$\mathbb{Z}$ 的任意非空子集若有上界，则有一个最大元。

(b) 若 $x \notin \mathbb{Z}$，证明恰好有一个 $n \in \mathbb{Z}$，使得 $n < x < n + 1$。

(c) 若 $x - y > 1$，证明至少有一个 $n \in \mathbb{Z}$，使得 $y < n < x$。

(d) 若 $y < x$，证明存在一个有理数 $z$，使得 $y < z < x$。
\end{exercise}

\begin{exercise}\label{exercise:ch01-整数与实数-10}
按照以下步骤证明任意正数 $a$ 恰好有一个正平方根：
\begin{enumerate}[(a)]
\item 证明，若 $x > 0$ 并且 $0 \leqslant h < 1$，则
\begin{align*}
(x + h)^{2} \leqslant x^{2} + h(2x + 1),
\end{align*}
\begin{align*}
(x - h)^{2} \geqslant x^{2} - h(2x).
\end{align*}
\item 设 $x > 0$。证明：若 $x^{2} < a$，则存在 $h > 0$ 使得 $(x + h)^{2} < a$；若 $x^{2} > a$，则存在 $h > 0$ 使得 $(x - h)^{2} > a$.
\item 给定 $a > 0$，设 $B$ 为使得 $x^{2} < a$ 的所有实数 $x$ 的集合。证明 $B$ 有上界并且至少含有一个正数。设 $b = \sup B$，证明 $b^{2} = a$.
\item 证明：若 $b$，$c$ 是一个正数并且 $b^{2} = c^{2}$，则 $b = c$.
\end{enumerate}
\end{exercise}

\begin{exercise}\label{exercise:ch01-整数与实数-11}
给定 $m \in \mathbb{Z}$，若 $m/2 \in \mathbb{Z}$，则称 $m$ 为偶数（even number），否则称 $m$ 为奇数（odd number）。
\begin{enumerate}[(a)]
\item 证明：若 $m$ 为奇数，则对于某一个 $n \in \mathbb{Z}$，有 $m = 2n + 1$. [提示：选取 $n$ 使得 $n < m/2 < n + 1$.]
\item 证明：若 $p$ 和 $q$ 都是奇数，则 $p \cdot q$ 是一个奇数，并且对于任何 $n \in \mathbb{Z}_{+}$，$p^{n}$ 也是一个奇数.
\item 证明：若 $a > 0$ 为有理数，则 $a = m/n$，其中 $m$，$n \in \mathbb{Z}_{+}$ 且不同时为偶数. [提示：取 $n$ 为 $\{x \mid x \in \mathbb{Z}_{+}$ 和 $x \cdot a \in \mathbb{Z}_{+}\}$ 的最小元.]
\item $\sqrt{2}$ 为无理数.
\end{enumerate}
\end{exercise}

\end{document}